\documentclass[10pt,a4paper]{amsart}
\usepackage{amsmath,amssymb,amsthm,graphicx}
\RequirePackage{amssymb,amsmath,amsthm,bm}
\RequirePackage{thmtools}
\RequirePackage[capitalize]{cleveref}
\providecommand{\tightlist}{%
  \setlength{\itemsep}{0pt}\setlength{\parskip}{0pt}}

%%%%% varphi instead of phi
\renewcommand{\phi}{\varphi}

\mathchardef\mhyph="2D 
\providecommand{\A}{\mathcal{A}}
\providecommand{\C}{\mathcal{C}}
\providecommand{\L}{\mathcal{L}}
\providecommand{\B}{\mathcal{B}}
\providecommand{\M}{\mathcal{M}}
\renewcommand{\S}{\mathcal{S}}
\providecommand{\ol}[1]{\overline{#1}}
\providecommand{\mbf}[1]{\mathbf{#1}}
\providecommand{\mf}[1]{\mathfrak{#1}}
\providecommand{\mrm}[1]{\mathrm{#1}}
\providecommand{\mc}[1]{\mathcal{#1}}
\providecommand{\ra}{\rightarrow}
\providecommand{\Ra}{\Rightarrow}
\providecommand{\La}{\Leftarrow}
\providecommand{\lr}{\leftrightarrow}
\providecommand{\hf}{\mathbb{HF}}
\providecommand{\spec}{DgSp}
\providecommand{\bSigma}{\bm\Sigma}
\providecommand{\bDelta}{\bm\Delta}
\providecommand{\bPi}{\bm\Pi}
\providecommand{\Pinf}[1]{\Pi^{\mathrm{in}}_{#1}}
\providecommand{\Sinf}[1]{\Sigma^{\mathrm{in}}_{#1}}
\providecommand{\dSinf}[1]{d\mhyph\Sigma^{\mathrm{in}}_{#1}}
\providecommand{\Dinf}[1]{\Delta^{\mathrm{in}}_{#1}}
\providecommand{\bp}{\bar p}
\providecommand{\bq}{\bar q}
\providecommand{\ba}{\bar a}
\providecommand{\bb}{\bar b}
\providecommand{\bc}{\bar c}
\providecommand{\bd}{\bar d}
\providecommand{\bx}{\bar x}
\providecommand{\by}{\bar y}
\providecommand{\bz}{\bar z}
\providecommand{\LO}{\mathrm{LO}}
\declaretheorem[name={Theorem},numberwithin=section]{theorem}
\declaretheorem[name={Proposition}, sibling=theorem]{proposition}
\declaretheorem[name={Lemma}, sibling=theorem]{lemma}
\declaretheorem[name={Fact}, sibling=theorem]{fact}
\declaretheorem[name={Claim}, numberwithin=theorem]{claim}
\declaretheorem[name={Definition},numberwithin=section,sibling=theorem, style=definition]{definition}
\declaretheorem[name={Remark}, numbered=no,style=definition]{remark}
\declaretheorem[name={Corollary}, sibling=theorem]{corollary}
\declaretheorem[name={Question},sibling=theorem]{question}
\newcommand{\T}{{\mathcal{T}}}
\newcommand{\SR}{SR}
\newcommand{\pSR}{pSR}
\newcommand{\Q}{\mathbb{Q}}
\newcommand{\Z}{\mathbb{Z}}
\newcommand{\K}{\mathcal{K}}
\newcommand{\N}{\mathcal{N}}
\renewcommand{\L}{\mathcal{L}}
\providecommand{\Ginf}[1]{\Gamma^{\mathrm{in}}_{#1}}
\title{A Borel reduction that eventually satisfies push-forward preserves Scott sentence complexity above a threshold}
\author{David Gonzalez, Dino Rossegger}
\date{arxiv:2305.07126v1; CC BY 4.0}
\begin{document}
\maketitle
\noindent\textit{Editorial note.} The counted statement is the author environment at source0.tex lines 696-704 and its proof environment opens at line 706, the line immediately after; the pairing is by adjacency and the retained proof range 706-718 is the whole environment.\par\medskip
\noindent\textit{Editorial note.} The retained preamble is the author style file dino.sty (package lines, the theorem declaration support, the notation and tuple macros and the theorem environments) followed by the ordinal and rank macros of the manuscript itself, reproduced verbatim; this is the same recipe as the previously published card of this article, so no package outside it is requested and extra\_packages stays empty.\par\medskip
\noindent\textit{Editorial note.} Metadata provenance: source\_id and source\_version are the values recorded in this source licence evidence file (sources/scott-linear-orders-2305.07126v1/LICENSE\_EVIDENCE.json: arxiv:2305.07126v1 / 2305.07126v1), which the card points at directly; the title and author list are transcribed from the authors\_index and title keys of the source record of the pinned copy. The card carries no counter setting and imposes no original numbering, since this article has no label-number dump and its numbering cannot be reconstructed from the retained ranges alone.\par\medskip
\noindent\textit{Editorial note.} No context span and no external reference are needed: the retained statement and proof use no label or citation from outside their own ranges.\par\medskip
\noindent\textit{Editorial scope.} One original result of the pinned author manuscript, counted once: preservation of Scott sentence complexity above a rank threshold by a Borel reduction that eventually satisfies push-forward (author label prop:push-forwardssc).\par\medskip
\noindent\textit{Allowed standard prerequisites (editorial).}\begin{itemize}
\item Borel reduction between classes of structures. A Borel reduction Phi from the models of sigma to the models of tau, measurable at some level of the Borel hierarchy, whose image is Borel.. It is the hypothesis of the proposition.
\item Push-forward property. Phi eventually satisfies push-forward when there is a rank delta such that every infinitary formula of rank at least delta over tau has a formula of the same rank over sigma pulling back along Phi.. It is the second hypothesis and the tool the proof uses.
\item Scott sentence complexity. SSC(A) is the least rank of a Scott sentence for A, and SR(A) is the Scott rank.. It is the quantity shown to be preserved above a threshold.
\end{itemize}\medskip
\begin{proposition}\label{prop:push-forwardssc}
   If $\Phi:Mod(\sigma)\to Mod(\tau)$ is a Borel reduction such that
   $[\Phi(X)]_{\cong}$ is Borel and that eventually satisfies push-forward, i.e.,
   there is $\delta$ such that for $\gamma\geq\delta$ and
   $\psi\in\Pinf{\gamma}$ there is $\psi^*\in\Pinf{\gamma}$ such that
   $$\A\models \psi \iff \Phi(A)\models\psi^*$$
   then there is an ordinal $\alpha$ such that $\text{SR}(\A)\geq\alpha$ implies
   that $SSC(\A)=SSC(\Phi(\A))$.
\end{proposition}

\begin{proof}
Say $\Phi$ is $\bSigma_\lambda^0$-measurable and the image of $\Phi$ is
$\bPi_\beta^0$. Let $\chi\in\Pinf{\beta}$ describe the image of $\Phi$ in
$Mod(\tau)$. 
Let $\A$ have a $\Gamma^{\mathrm{in}}_\epsilon$ Scott sentence for
$\epsilon>\delta$, say $\phi$. Then $\phi^*\land \chi$ is a
$\Gamma^\mathrm{in}_\epsilon$ Scott sentence for $\Phi(\A)$. On the other hand
say $\Phi(\A)$ has a $\Gamma^{\mathrm{in}}_\epsilon$ Scott sentence for any
$\epsilon$, then by the pullback theorem $\A$ has a
$\Gamma^{\mathrm{in}}_{\lambda+\epsilon}$ Scott sentence. So for any $\A$ with
$\SR(\A)>\max(\lambda\cdot \epsilon, \delta,\beta)$, Scott sentence complexity
is preserved.
\end{proof}

\end{document}
